\worksheettitle
  {M03-R. Сохраняем длину}
  {\modulemark{M03-R} Коррекционная карточка}

\studentnameblock

\begin{warningbox}[Назовите ошибку]
  Нельзя менять запись числа, не назвав соотношение единиц. Сначала выясните,
  сколько требуемых единиц содержится в одной исходной единице.
\end{warningbox}

\begin{practicebox}[Шаг 1. Восстановите опору]
  \[
    1\text{ см}=\answerblank[12mm]\text{ мм},\quad
    1\text{ м}=\answerblank[15mm]\text{ см},\quad
    1\text{ км}=\answerblank[18mm]\text{ м}.
  \]
\end{practicebox}

\begin{practicebox}[Шаг 2. Объясните направление]
  Переводим $4$ м в сантиметры. Сантиметр меньше метра, поэтому число должно
  \answerblank[28mm].
  \[
    4\text{ м}=4\cdot\answerblank[18mm]\text{ см}
      =\answerblank[25mm]\text{ см}.
  \]
\end{practicebox}

\begin{practicebox}[Шаг 3. Составная длина]
  \[
    3\text{ м }7\text{ см}=3\cdot\answerblank[18mm]\text{ см}
      +7\text{ см}=\answerblank[26mm]\text{ см}.
  \]
\end{practicebox}

\begin{practicebox}[Шаг 4. Исправьте]
  Ученик написал: $5\text{ км }20\text{ м}=520\text{ м}$.

  Причина: \answerblank[95mm]

  Исправление: \answerblank[95mm]
\end{practicebox}

\begin{checkpointbox}[Повторная проверка]
  \begin{enumerate}[label=\textbf{\arabic*.}]
    \item $2\text{ м }35\text{ см}=\answerblank[28mm]\text{ см}$.
    \item $3040\text{ м}=\answerblank[18mm]\text{ км }
      \answerblank[22mm]\text{ м}$.
    \item Объясните учителю, почему ответы сохраняют исходную длину.
  \end{enumerate}
\end{checkpointbox}
